\documentclass[10pt,a4paper]{amsart}
\usepackage[margin=19mm]{geometry}
\usepackage{amsmath,amssymb,mathtools}
\usepackage[scr=rsfs,cal=euler]{mathalfa}
\usepackage{xfrac}
\usepackage[unicode,hidelinks,bookmarks=false]{hyperref}
\newcommand{\ie}{i.e.\ }
\newcommand{\cf}{cf.\ }
\newcommand{\resp}{respectively\ }
\newcommand{\Subsection}{\subsection}
\providecommand{\nobreakdash}{\nobreak\hbox{-}}
\setlength{\emergencystretch}{2em}
\multlinegap0pt
\usepackage{amssymb}
\usepackage{mathtools}
\usepackage{xcolor}
\newtheorem{proposition}{Proposition}
\newtheorem{theorem}{Theorem}
\newtheorem{lemma}{Lemma}
\newcommand{\CC}{\mathcal{C}}
\newcommand{\On}[1]{\operatorname{O}(#1)}
\newcommand{\PP}{\mathcal{P}}
\newcommand{\R}{\mathbb{R}}
\newcommand{\SO}{\operatorname{SO}}
\newcommand{\SOn}[1]{\operatorname{SO}(#1)}
\newcommand{\lin}{\operatorname{lin}}
\newcommand{\sphere}[1]{\mathbb{S}^{#1}}
\title{Monotone invariant valuations on convex cones}
\author{Martin Lotz}
\date{Source: arXiv:2609.09335v2 (CC BY 4.0)}
\begin{document}
\maketitle

\noindent\emph{Source and licence.} Monotone invariant valuations on convex cones. Version arXiv:2609.09335v2 (CC BY 4.0), distributed by arXiv under the Creative Commons Attribution 4.0 International licence, http://creativecommons.org/licenses/by/4.0/, as recorded in the arXiv licence metadata for this exact version and shown on the abstract page https://arxiv.org/abs/2609.09335v2. Full licence notice: \texttt{licenses/arxiv-2609.09335-CC-BY-4.0.txt}.\par\smallskip

\noindent\textit{Editorial scope.} Counted unit: the article's theorem thm:regularity (source0.tex lines 823--830). The article's complete original proof of it is delivered with it (source0.tex lines 832--885, one \texttt{proof} environment of 54 lines). No mathematics is written, repaired or re-derived: the statement and the proof are the article's own lines, in the article's own notation, and no retained line is substituted or re-flowed.  Retained uncounted as the source-local context the proof needs: lines 351--354; lines 366--375; lines 676--689; lines 727--732; lines 797--803.  Nothing was omitted from any retained span, so the package carries no omission record.  No retained line contains a citation, so the wrapper carries no bibliography and no bibliography entry.  No retained line contains a figure environment, an image-inclusion command or drawing code, so no image is delivered and the package declares no asset dependency.  Editorial formatting only, in addition: an AMS \texttt{amsart} preamble that restates the article's own declarations and macros used by the retained lines, namely the environments \texttt{proposition}, \texttt{theorem}, \texttt{lemma} on separate counters, together with the standard packages those lines need.  The retained environments print in this document as Proposition 1, Theorem 1, Lemma 1 in order of appearance, whereas the article prints the same items with the numbering of its own full document.  The complete native source of the article as distributed in the arXiv e-print for this exact version is bundled unchanged as \texttt{sources/209802/source0.tex}.
\begin{equation}\label{eq:orthant-triangular}
  U_j(O_k)=0\quad(j\ge k),\qquad
  U_{k-1}(O_k)=v_k(O_k)=2^{-k}.
\end{equation}

\begin{proposition}\label{prop:polyhedral-sandwich}
For every $C\in\CC_p^*(\R^d)$ there are sequences
$P_r,Q_r\in\PP_p^*(\R^d)$ such that
\begin{equation*}
  P_r\subseteq C\subseteq Q_r,
  \qquad P_r\rightarrow C,
  \qquad Q_r\rightarrow C
\end{equation*}
in the spherical Hausdorff topology.
\end{proposition}

\begin{theorem}\label{thm:simple-integrable}
Let $\eta$ be a simple $\SOn d$-invariant valuation on
$\PP_p^*(\R^d)$.  Suppose that for every pointed simplicial cone $F$
of dimension $d-1$, the function
\[
  x\longmapsto\eta(F+\R_{\ge0}x)
\]
is integrable on $\sphere{d-1}\setminus\lin F$.
Then
\[
  \eta(P)=\bar\eta(\R^d)v_d(P)
  \qquad\text{for all }P\in\PP_p^*(\R^d).
\]
\end{theorem}

\begin{lemma}\label{lem:normalization}
Let $\mu$ be a monotone $\SOn d$-invariant valuation on
$\CC_p^*(\R^d)$, and let $\alpha$ be its common value on rays.
Then $\mu_0:=\mu-2\alpha U_0$ is monotone, invariant, nonnegative,
bounded, and vanishes on rays.
\end{lemma}

\begin{lemma}\label{lem:measurable-coning}
Let $H\subset\R^d$ be a linear hyperplane and let $\mu$ be a monotone
valuation on $\CC_p^*(\R^d)$.
If $C\in\CC_p^*(H)$ is full-dimensional in $H$, then
$x\mapsto\mu(x*C)$ is measurable on $\sphere{d-1}\setminus H$
for completed spherical measure.
\end{lemma}

\begin{theorem}\label{thm:regularity}
Every monotone $\SOn d$-invariant valuation on $\CC_p^*(\R^d)$
has a unique representation
\begin{equation}\label{eq:pointed-representation}
  \mu=\sum_{j=0}^{d-1}a_jU_j.
\end{equation}
In particular, it is continuous and $\On d$-invariant.
\end{theorem}

\begin{proof}
Uniqueness follows from \eqref{eq:orthant-triangular}.
We prove existence by induction on $d$, normalising $\mu$ by
Lemma~\ref{lem:normalization} to be bounded, nonnegative, and zero
on rays.

For $d=2$, set $\nu=0$: the normalised valuation already vanishes
on all lower-dimen\-sional pointed cones.
For $d\ge3$, assume the representation in dimension $d-1$ and
construct $\nu$ as follows.
For a hyperplane $W$, the restriction of $\mu$ to $\CC_p^*(W)$ is
monotone and $\SO(W)$-invariant, since rotations of $W$ extend to
rotations of $\R^d$ fixing its normal line pointwise.
By induction,
\[
  \mu|_{\CC_p^*(W)}=\sum_{j=1}^{d-2}b_jU_j^W.
\]
The constant coefficient is zero because $\mu$ vanishes on rays.
Uniqueness and rotation invariance make the coefficients independent
of $W$.  Intrinsicness of the Grassmann angles shows that
\[
  \nu:=\sum_{j=1}^{d-2}b_jU_j
\]
agrees with $\mu$ on every lower-dimensional nonzero pointed cone.

In either case, $\eta:=\mu-\nu$ is simple, bounded, and
$\SOn d$-invariant.
Here boundedness of $\nu$ follows directly from its intrinsic-volume
representation.  Neither monotonicity nor nonnegativity of $\eta$
is assumed.

For every pointed simplicial facet $F$,
Lemma~\ref{lem:measurable-coning} makes $x\mapsto\mu(x*F)$ measurable.
The function $x\mapsto\nu(x*F)$ is continuous off $\lin F$.
Thus the apex functions for $\eta$ are measurable and bounded.
Theorem~\ref{thm:simple-integrable} gives a constant $c$ such that
\[
  \mu(P)=g(P),\qquad
  g:=\nu+cv_d=\nu+cU_{d-1},
  \qquad P\in\PP_p^*(\R^d).
\]

For $C\in\CC_p^*(\R^d)$, take the approximations from
Proposition~\ref{prop:polyhedral-sandwich}.  Monotonicity of the
original $\mu$ gives
\begin{equation}\label{eq:monotone-sandwich}
  g(P_r)=\mu(P_r)\le\mu(C)\le\mu(Q_r)=g(Q_r).
\end{equation}
Both bounds converge to $g(C)$ by continuity of the known function $g$.
Hence $\mu(C)=g(C)$.  Restoring the constant removed in the
normalisation proves \eqref{eq:pointed-representation} and completes
the induction.  Continuity and orthogonal invariance follow from
the corresponding properties of the $U_j$.
\end{proof}

\end{document}
